\begin{frame}[t]{Forecast electricity demand}
\lead{GEFCom2012: hourly demand from a US utility}
\begin{center}\resizebox{.94\linewidth}{!}{\fig{case-load-pipeline}}\end{center}
\small
\begin{columns}[T,onlytextwidth]
\begin{column}{.55\textwidth}
20 zones; 11 weather stations; 4.5 years of records.\par
\uncover<2->{\begin{definitionbox}Team TinTin: boost \textbf{smooth spline functions}.\end{definitionbox}}
\end{column}
\begin{column}{.41\textwidth}
\uncover<2->{\begin{takeaway}\textbf{5th of 105 teams}\\Reported competition ranking.\end{takeaway}}
\end{column}
\end{columns}
\vfill{\scriptsize\href{https://robjhyndman.com/papers/kaggle-competition.pdf}{Ben Taieb \& Hyndman (2014), Abstract; Sections 3--4.}}
\note{The study used separate hourly models and log-transformed demand. Spline learners extend the same residual-fitting principle beyond trees. This is a competition result on real utility records, not evidence of operational deployment or a universal ranking of algorithms.}
\end{frame}

\begin{frame}[t]{What is known when we predict?}
\lead{Eight historical weeks to reconstruct; one future week to forecast}
\begin{center}\resizebox{.86\linewidth}{!}{\fig{case-load-timeline}}\end{center}
\vfill{\scriptsize\href{https://robjhyndman.com/papers/kaggle-competition.pdf}{Ben Taieb \& Hyndman (2014), Section 3. Timeline is schematic.}}

\end{frame}
